\chapter{Programming Model}

This chapter summarizes the kernel and userspace rules for the asynchronous,
shard-per-core model.

\section{Invariants}

The design relies on these rules:

\subsection{Only the owning LP mutates its state}

Per-LP data is accessed only from the LP that owns it. Cross-LP mutation goes
through message dispatch (\texttt{ShardMailbox<M>}, IPI closure, or IPI RPC),
never through a shared lock acquired from a foreign LP.

\subsection{Messages are owned values}

Everything that crosses a shard boundary is an owned value --- a typed
\texttt{M} in a \texttt{ShardSender<M>}, an IPI closure, or a memory object.
No shared references, no \texttt{Arc<Mutex<...>{>}} as the normal programming
model. \texttt{Send + 'static} boundaries enforce this at compile time.

\subsection{References to shard-local state must not escape}

\texttt{ShardLocal<T>::with(fn)} provides \texttt{\&mut T} inside a closure.
The closure must not store the reference, send it to another thread, or cross
an \texttt{.await} point. The borrow is verified at runtime (owner-check plus
re-entrancy flag).

\subsection{Work stays on its assigned LP}

A thread spawned on LP $N$ stays on LP $N$. To submit work to another LP, use
\texttt{try\_run\_on\_lp(target, closure)} or \texttt{ShardSender::try\_send}.

\subsection{Observability returns owned snapshots}

Never expose borrowed references into runtime internals. Diagnostic and
introspection code captures state by value. Self-tests and the demo illustrate
this pattern.

\section{Kernel-side primitives}

\subsection{\texttt{PerLp<T>}}

Lock-wrapped per-LP container. Use for data accessed from interrupt handlers,
IPI handlers, or multiple LPs:

\begin{verbatim}
// Declaration:
static SCHEDULER: PerLp<RoundRobin> = PerLp::new();

// Access:
SCHEDULER.with(|scheduler| {
    scheduler.submit_ready_thread(tid);
});
\end{verbatim}

The lock masks interrupts during the critical section on architectures that
require it.

\subsection{\texttt{ShardLocal<T>}}

Lock-free per-LP container for single-LP, thread-only data:

\begin{verbatim}
// Declaration:
static MAILBOX: ShardLocal<ShardMailbox<Command>> = ShardLocal::new();

// Access (only from owning LP, not from interrupt context):
MAILBOX.with(|mb| {
    mb.try_send(Command::Ping)?;
});
\end{verbatim}

Panics on wrong-LP access or re-entrant access.

\subsection{\texttt{ShardMailbox<M>}}

A typed, bounded per-LP queue has cloneable \texttt{ShardSender<M>} producers and
one \texttt{ShardReceiver<M>} consumer per LP, accessed through
\texttt{ShardLocal}:

\begin{verbatim}
let sender: ShardSender<Command> = mailbox.sender();
// Returns Err(data) if full.
sender.try_send(Command::Process { data })?;
\end{verbatim}

\subsection{IPI closures}

\texttt{try\_run\_on\_lp(target\_lp, || \{ ... \})}. Boxes arbitrary work and
delivers it to the target LP's IPI queue. Returns the closure on full queue
(backpressure). Use for kernel-level cross-LP work where typed mailboxes are not
appropriate (e.g., TLB shootdown, scheduler commands).

\subsection{Completion API}

The API provides \texttt{submit}, \texttt{submit\_worker}, \texttt{complete},
\texttt{poll}, \texttt{wait}, \texttt{cancel}, and \texttt{close}. The exit-observer path
(\texttt{submit\_worker}) is the intended default for fire-and-forget async
kernel work:

\begin{verbatim}
let cap = submit_worker(|| {
    // Worker thread: do async work, then exit.
    // Completion fires automatically on exit.
});
let result = wait(cap)?; // Block until worker exits.
\end{verbatim}

\subsection{CQ ring}

\texttt{CompletionQueueRing} for zero-syscall completion draining.
Kernel side: \texttt{write(cap, result)}. Consumer side:
\texttt{read() -> Option<CqEntry>}.

\section{Userspace programming model}

\subsection{The launch contract}

A CharlotteOS userspace program is a \texttt{no\_std} ELF binary. It uses
the \texttt{catten-rt} crate (the CRT, or C runtime) and the \texttt{entry!}
macro:

\begin{verbatim}
#![no_std]
#![no_main]

use catten_rt::entry;

entry!(main);

fn main(ctx: Context) -> ! {
    // ctx provides:
    // - A connection to the name service
    // - The default CQ handle
    // - Heap configuration
    // - Bootstrap capabilities (device grants, etc.)

    let name_service = ctx.name_service();
    let nic_driver = name_service.lookup("driver.nic.v1")?;

    // ... service logic ...

    loop {
        // Poll tasks, drain CQ, wait for events.
    }
}
\end{verbatim}

The program does not define \texttt{\_start}, a panic handler, or an allocator.
These are provided by \texttt{catten-rt}. The \texttt{entry!} macro wires them
together, validates the launch ABI version, sets up the heap, and calls
\texttt{main(ctx)}. A panic invokes the domain-abort syscall: every thread in
the address space is retired, so no sibling can continue after shared Rust
state or a userspace lock has been left inconsistent.

\subsection{Memory ownership in Rust}

The register-level syscall ABI deliberately exposes capability handles as
integers. Application code should use the linear wrappers in
\texttt{catten\_rt::owned}:

\begin{itemize}
    \item \texttt{OwnedMemory} closes its capability on drop and is neither
        \texttt{Copy} nor \texttt{Clone}.
    \item \texttt{MappedMemory<ReadOnly>} and
        \texttt{MappedMemory<Writable>} own both the mapping and capability;
        only the writable state can produce \texttt{\&mut [u8]}.
    \item \texttt{DmaTransfer} is created only from unmapped
        \texttt{OwnedMemory} and returns CPU ownership after synchronous DMA
        unmapping.
    \item \texttt{ReadOperation<'a>} retains its
        \texttt{\&'a mut [u8]} until completion, including the
        cancel-and-wait path in its destructor.
\end{itemize}

The raw \texttt{catten-syscall} interface remains available to runtimes and
drivers. Coherent \texttt{dma\_map} is explicitly unsafe because the driver
must synchronize CPU references and device access itself. Mixing raw
operations with an adopted ownership wrapper violates the wrapper's safety
contract.

\subsection{Bootstrap authority}

The \texttt{Context} carries the domain's bootstrap capabilities:

\begin{itemize}
    \item \texttt{name\_service()} --- a connection to the name service, for
        looking up other services.
    \item \texttt{default\_cq()} --- the CQ handle for submission and waiting.
    \item \texttt{device\_grants()} --- any \texttt{MmioRegion} or
        \texttt{Interrupt} capabilities granted to this domain.
    \item \texttt{heap\_config()} --- memory available for the allocator.
\end{itemize}

The \texttt{Context} is the only way a userspace program acquires capabilities
at startup. There is no ambient authority, no global filesystem, no
\texttt{/dev} to explore.

\subsection{Caller identity is kernel authority}

Syscalls derive the caller's address space from trusted thread and trap context.
Userspace cannot select another domain's address space or capability table.

\subsection{Backpressure is normal}

Bounded queues and capability tables return \texttt{WouldBlock}. The correct
response is to retry, yield, or propagate backpressure upstream. Do not assume 
an unbounded queue.

\subsection{Sitas integration}

When writing code that uses sitas on CharlotteOS:

\begin{itemize}
    \item A sitas shard is an LP-affine CharlotteOS user thread.
    \item \texttt{CharlotteReactor::new(lp\_id)} binds a reactor to the calling
        shard's LP.
    \item Async operations return kernel-owned completion handles. A completion
        handle is opaque: wait on it, poll it, close it, or pass it through a
        typed channel.
    \item Cross-shard work is an owned message. Prefer sitas channels and
        futures over shared mutable state.
    \item The executor parks on \texttt{CQ\_WAIT} (via \texttt{ShardParker}),
        not on a busy-loop.
\end{itemize}

\section{Rule-of-thumb guidance}

\begin{itemize}
    \item \textbf{Choose the right container:} \texttt{PerLp<T>} for
        interrupt-accessed or cross-LP data; \texttt{ShardLocal<T>} for
        single-LP, thread-only data.
    \item \textbf{Choose the right channel:} \texttt{ShardMailbox<M>} for typed
        cross-LP communication within the kernel. \texttt{ShardSender/Receiver}
        for typed cross-shard communication within userspace. Endpoint IPC
        for cross-domain communication.
    \item \textbf{Use \texttt{submit\_worker} for fire-and-forget async work.}
        When the worker returns, the exit-observer completes the capability.
    \item \textbf{Register a CQ ring at address-space creation} so completions
        are visible to userspace without per-entry syscalls.
    \item \textbf{Accept backpressure.} \texttt{WouldBlock} or \texttt{Err(m)}
        means the receiver is slow, not that the system is broken. Retry or
        propagate upstream.
    \item \textbf{Buffers have one owner.} While an async operation owns a
        buffer, userspace must not mutate or free it. Cancellation still ends
        in a terminal completion.
    \item \textbf{Never accept ASID or page-table roots from userspace.} Caller
        identity comes from the trap frame. This is the security invariant.
    \item \textbf{Do not smuggle shared kernel state into userspace.} Shared
        pages (CQ ring, config page) are data queues with narrow layouts, not
        references to kernel objects.
\end{itemize}

\section{The syscall ABI}

The userspace--kernel boundary exposes these operations:

\begin{description}
    \item[\texttt{submit(op\_code, buffer) -> OperationId}] Start an async
        operation. Returns immediately. Buffer ownership transfers to the kernel.

    \item[\texttt{wait(cq\_id) -> Vec<CompletionEntry>}] Block until a
        completion arrives on the given CQ. Returns all pending entries.

    \item[\texttt{poll(op\_id) -> Option<Completed>}] Non-blocking check.
        Returns \texttt{None} while in flight.

    \item[\texttt{cancel(op\_id) -> CancelState}] Request cancellation. May
        complete asynchronously.

    \item[\texttt{close(op\_id)}] Release an operation ID after its terminal
        completion was observed.
\end{description}

For IPC:

\begin{description}
    \item[\texttt{endpoint\_create(interface, capacity) -> EndpointCap}]
    \item[\texttt{connection\_mint(endpoint, rights) -> ConnectionCap}]
    \item[\texttt{scalar\_send(connection, opcode, arg)}]
    \item[\texttt{scalar\_call(connection, opcode, arg) -> ReplyToken}]
    \item[\texttt{memory\_send(connection, opcode, attachments)}]
    \item[\texttt{reply(token, result)}]
    \item[\texttt{receive(endpoint) -> Message}]
\end{description}

The full ABI is described in
\repofile|docs/architecture/async-syscall-abi.md| and implemented in
\repofile|crates/catten/src/syscall/mod.rs|.

\section{Why the programming model matters for critical software}

\begin{enumerate}
    \item \textbf{Invariants have enforcement points.} \texttt{ShardLocal<T>}
        rejects wrong-LP and re-entrant access; \texttt{ShardSender<M>} moves
        values across the typed boundary. Unsafe code and incorrect primitive use
        remain review obligations.

    \item \textbf{The launch contract is minimal.} A userspace program receives
        exactly what it needs and nothing more. There is no global state, no
        inherited file descriptors, no ambient filesystem. The bootstrap
        capabilities are the program's entire universe of authority.

    \item \textbf{Messages, not shared memory, are the default.} The safe,
        preferred path is to send an owned value through a bounded channel.
        Shared memory with locks is the exception and not the rule. This reduces
        the number of places where concurrency bugs can hide.

    \item \textbf{The ABI is auditable.} Userspace wrappers and the kernel
        dispatcher share the typed \texttt{SyscallNumber} enum, and dispatch is
        exhaustive. The source enum, rather than a duplicated count here, defines
        the current surface.

    \item \textbf{The programming model is testable.} Self-tests exercise
        owner checks, re-entrancy guards, selected backpressure paths, and
        cancellation behavior. Coverage is substantial but not exhaustive.
\end{enumerate}

\endinput
